% META: {"id":"TCOMPL-BAYES-CO-003","chapitre":"TCOMPL-INFERENCE-BAYESIENNE","type_objet":"corrige","exercice_ref":"TCOMPL-BAYES-EX-003","capacites_codes":["C4"],"sources_inspiration":[],"status":"approved","origin":"generated","status_history":["generated","audited","accepted","approved"],"release_acceptance":"RELEASE_OWNER_BATCH_ACCEPTANCE","acceptance_closure_digest":"sha256:9b3ccf9a81c5520fbb7e03b7b2d3e7bf2057a02834b6d908bf3be5f49f3b553f","acceptance_receipt_digest":"sha256:4fad86f0282e0fa4e0419cca1ad6379ae7af5a280c04a4a90880f90a1c044242"}
% BEGIN-VERIFY
% from sympy import Rational
% Se = Rational(90, 100)
% Sp = Rational(95, 100)
% p = Rational(2, 100)
% PB = Se*p + (1-Sp)*(1-p)
% assert PB == Rational(67, 1000)
% VPP = Se*p / PB
% assert VPP == Rational(18, 67)
% END-VERIFY
\begin{corrige}{TCOMPL-BAYES-EX-003}
\textbf{1.} $P(T) = Se\times p + (1-Sp)\times(1-p) = 0{,}90\times0{,}02 + 0{,}05\times0{,}98 = 0{,}018+0{,}049 = 0{,}067$.

\textbf{2.} $VPP = \dfrac{Se\times p}{P(T)} = \dfrac{0{,}018}{0{,}067} = \dfrac{18}{67} \approx 0{,}269$.

Malgré un test de bonne qualité ($90\%$ de sensibilité, $95\%$ de spécificité), seulement $26{,}9\%$ des personnes testées positives sont réellement malades, en raison de la faible prévalence ($2\%$).
\end{corrige}
